% concatenated from ICSF-hyperbolic-plane-for-arxiv.tex
\documentclass[12pt]{amsart}
%\usepackage{extsizes}
 

\usepackage{amssymb, amsmath, amsthm, amsfonts}
\usepackage{hyperref}
\usepackage{xcolor,graphicx}



%%%%%%
 \setlength{\textwidth}{15cm}
\setlength{\oddsidemargin}{1cm}
\setlength{\evensidemargin}{1cm}
\setlength{\textheight}{21cm}
\setlength{\parskip}{2mm}
\setlength{\parindent}{0em}
\setlength{\headsep}{1.5cm}
%%%%%

\theoremstyle{plain}
\newtheorem{theorem}{Theorem}[section]

\newtheorem{lemma}[theorem]{Lemma}
\newtheorem{corollary}[theorem]{Corollary}
\newtheorem{proposition}[theorem]{Proposition}
\newtheorem{conjecture}{Conjecture}

 


\theoremstyle{definition}
\newtheorem{definition}[theorem]{Definition}
\newtheorem{remark}[theorem]{Remark}
\newtheorem{example}[theorem]{Example}
\newcommand{\N}{\mathbb{N}}
\newcommand{\Z}{\mathbb{Z}}
\newcommand{\Q}{\mathbb{Q}}
%\renewcommand{\H}{\mathbb{H}}
\def\h{\mathbb H}
\def\r{\mathbb R}
\def\c{\mathbb C}
%\renewcommand{\tanh}{\tanh} %renewcommand jer \tanh već postoji
\newcommand{\ch}{\text{ch}}
\newcommand{\sh}{\text{sh}}
\newcommand{\arch}{\text{arch}}
\newcommand{\arctg}{\text{arctg}}

%\usepackage{showkeys}

%-----------------------------------------------------------------------------------------------------

\begin{document}
\title{The inverse curve shortening flow on the hyperbolic plane }
\author{Ivan Krznari\'{c}}
\address{Faculty of Economics and Business. University of Zagreb}
\email{ikrznaric@efzg.hr}
\author{Rafael L\'opez}
\address{Departamento de Geometr\'{\i}a y Topolog\'{\i}a Universidad de Granada 18071 Granada, Spain}
\email{rcamino@ugr.es}

\subjclass[2020]{35J60,   53E10, 53C21, 53C44, 53C50, 58J05.}  
\keywords{inverse curvature flow, soliton, hyperbolic plane, parabolic, conformal}

\begin{abstract}
We study the inverse curve shortening flow in the hyperbolic plane $\h^2$. We classify all solitons with respect to parabolic  and conformal vector fields of $\h^2$. In the upper half-plane model of $\h^2$, we prove that parabolic solitons are all graphs on the $y$-axis, whereas  conformal solitons are graphs on the $x$-axis. We study the concavity of these solitons and when they approach the coordinate axes.
\end{abstract}

\maketitle

%%%%%%%%%%%%%%%%%%%%%%%%%%%
\section{Introduction}

Let $\gamma\colon I\to\h^2$ be a curve in the hyperbolic plane $\h^2$, parametrized by arc length. If $J\subset\r$ is an interval with $0\in J$, a one-parameter family of parametrized curves $\{\gamma_t\colon I\to\h^2, t\in J\}$ evolves by the  inverse curve shortening flow (for abbreviation, ICSF) with initial condition $\gamma$,   if 
\begin{equation}\label{flow}
\left\{
\begin{split}
\left(\frac{\partial \gamma_t}{\partial t}\right)^\perp&= - \frac{1}{\kappa(t,\cdot )} N(t,\cdot ), \\ 
\gamma_0&=\gamma. \end{split}\right.
\end{equation}
 where $\kappa(t,\cdot)$ and $N(t,\cdot)$ denote the curvature and the unit normal vector of $\gamma_t$, respectively, and $ \left(\frac{\partial \gamma_t}{\partial t}\right)^\perp$ is the normal component of the velocity $ \gamma_t$. In arbitrary dimension, the inverse mean curvature flow (IMCF) in Euclidean space $\r^{n+1}$ was first studied by Huisken, who proved that compact, mean convex, star-shaped hypersurfaces converge to round spheres under the flow \cite{ge}: see also \cite{cd,ch}. Non-compact solitons of the IMCF are known to exist \cite{dhw,hi}. 
The IMCF  has been considered when the ambient space is    the hyperbolic space: \cite{bda,ge2,ge3,hw}, and it has been   extended to other ambient spaces \cite{cl,kp,lw,lu,ns,pi1,pi2,sc,zh}

In this paper, we aim to provide additional results about non-compact solutions to the ICSF in the hyperbolic plane $\h^2$ by studying solitons with respect to parabolic and conformal vector fields. After introducing some basic facts of the upper half-plane model of $\h^2$ (Section \ref{sec2}), we present explicit examples of the ICSF for curves of constant curvature (Section \ref{sec3}).  The last two sections of the paper are devoted to the study of solitons of the ICSF with respect to  specific vector fields. Consider the upper half-plane model of $\h^2$ with standard coordinates $(x,y)$. In Section \ref{sec5}, we classify the  solitons of the ICSF corresponding to the parabolic vector field $\partial_x$. These solitons are all graphs on the $y$-axis; some of them approach the ideal boundary of $\h^2$ orthogonally, while others have horizontal asymptotes. In Section \ref{sec6}, we classify solitons of the ICSF corresponding to the conformal vector field $\partial_y$.  These solitons are concave graphs on bounded intervals on the $x$-axis with infinite slope at the end points. 
 
 

%%%%%%%%%%%%%%%%%%%%%%%
\section{Preliminaries} \label{sec2}
%%%%%%
In this paper, we work with the upper half-plane model of the hyperbolic plane $\h^2$, namely, 
$$\r^2_+= \left\{ (x,y) \in \r^2\colon y > 0\right\}$$
 endowed with the Riemannian metric 
 $$\langle,\rangle = \frac{dx^2 + dy^2}{y^2}.$$
 The ideal boundary $\h^2_\infty=\{y=0\}\cup\{\infty\}$ is the one-point compactification of the horizontal line $y=0$.  


Let $\gamma\colon I\to\r^2_+$ be a regular curve, parametrized as $\gamma(s)=(x(s),y(s))$. There is a well-known relationship between 
the curvature $\kappa$ of the curve $\gamma$ in $\r^2_+$ with the hyperbolic metric, and its Euclidean curvature $\kappa_e$, since the hyperbolic metric is conformal to the Euclidean one. The unit normal vector $N$ and the hyperbolic curvature $\kappa$ of $\gamma$ are given by 
\begin{eqnarray}
N&=&y\frac{(-y',x')}{\sqrt{x'^2+y'^2}},\label{ke1}\\
\kappa&=&y\kappa_e+\frac{x'}{\sqrt{x'^2+y'^2}},\qquad \kappa_e=\frac{x'y''-x''y'}{(x'^2+y'^2)^{3/2}}.\label{ke2}
\end{eqnarray} 
As a consequence, in the upper half-plane model of $\h^2$, curves of constant hyperbolic curvature coincide with the curves of constant Euclidean curvature in $\r^2_+$.  Thus, they are precisely the intersections of straight lines and Euclidean circles with $\r^2_{+}$. The classification of these curves according to the curvature $\kappa$ is the following:  
\begin{enumerate}
\item Geodesics ($\kappa=0$). These are vertical lines and half-circles orthogonal to the line $y=0$.  
\item Equidistant lines ($0<\kappa<1$). These are non-geodesics arcs of circles supported on the line $y=0$ and non-vertical, non-horizontal straight lines.  
\item Horocycles ($\kappa=1$). These are circles tangent to the line $y=0$ or horizontal lines.   
\item Hyperbolic circles  ($\kappa>1$). These are Euclidean circles entirely contained in $\r^2_+$.  
\end{enumerate}

We now introduce the notion of a soliton of the ICSF by a Killing vector field. Let $X\in\mathfrak{X}(\h^2)$ be a Killing vector field. Then there is a one-parameter group of isometries $\Phi\colon\h^2\times\r\to\h^2$, $\Phi=\Phi(p,t)$,  such that 
\begin{equation}\label{ft}
\frac{d}{dt}\Phi(p,t)=X(\Phi(p,t)),\quad \mbox{and}\quad \Phi(p, 0)=p,
\end{equation}
for all $p\in\h^2$ and $t\in\r$.

Let $\{\gamma_t\colon I\to\h^2\colon t\in J\}$ be  a solution of the ICSF \eqref{flow}. We say that $\gamma=\gamma_0$ that is \textit{a soliton} with respect to   $X$ if $\Phi^{-1}\circ\gamma_t$ is stationary in the normal direction. Since each $\Phi_t$ is an isometry of $\h^2$, we have $N(t,s)=\Phi(N(s),t)$, where $N(s)$ denotes the unit normal vector of $\gamma_0$. Therefore,

\begin{equation*}
\begin{split}
\frac{1}{\kappa(t,s)}&=- \langle \frac{\partial \gamma_t}{\partial t}, N(t,s)\rangle \\
		       & = -\langle \frac{d}{dt}\Phi(\gamma_0,t),  \Phi(N(s),t)\rangle \\
		       & =- \langle X(\gamma_t), \Phi(N(s),t) \rangle.
\end{split}
\end{equation*}
Since this identity holds for all $t\in J$, setting $t=0$ gives 
$$\frac{1}{\kappa(s)}=-\langle X(\gamma_0),N(s)\rangle.$$
We have therefore proved the following characterization of solitons of the ICSF with respect to a Killing vector field, which we will be using throughout the paper.

\begin{proposition}
Let $X\in\mathfrak{X}(\h^2)$ be a Killing vector field. A curve $\gamma$ is a soliton of the ICSF with respect to $X$ if and only if 
\begin{equation}\label{soliton}
\frac{1}{\kappa}=-\langle N,X\circ\gamma\rangle,
\end{equation}
where $\kappa$ and $N$ are the curvature and the unit normal vector of $\gamma$, respectively.
\end{proposition}

The space of Killing vector fields on $\h^2$ is three-dimensional. A basis is given by 
\begin{equation}\label{kv}
\begin{split}
 X_1(x,y)&=\partial_x,\\
 X_2(x,y)&=x\partial_x+y\partial_y,\\
 X_3(x,y)&=(y^2-x^2)\partial_x+2xy\partial_y,
 \end{split}
\end{equation}
where $\{\partial_x,\partial_y \}$ denotes the canonical basis of  $\mathfrak{X}(\h^2)$. These vector fields correspond to one-parameter subgroups of isometries of $\h^2$. In complex notation $z=x+iy\in\r^2_+\subset\mathbb{C}$, and taking into account \eqref{ft}, the description of these subgroups is the following.
\begin{enumerate}
\item Parabolic translations. $\mathcal{P}=\{P_t\colon t\in\r\}$, where $P_t(z) = z + t$. Then 
$$X_1(z)=\frac{d}{dt}{\Big|}_{t=0}P_t(z).$$
\item Hyperbolic translations. $\mathcal{H}=\{H_t\colon t\in\r\}$, where $H_t(z) = e^{t}z$. Then 
$$X_2(z)=\frac{d}{dt}{\Big|}_{t=0}H_t(z).$$ 
\item Rotations. $\mathcal{R}=\{ R_t\colon t\in\r\}$, where $R_t(z) = \frac{\cos(t)z - \sin(t)}{\sin(t)z +\cos(t)}$. Then $$X_3(z)=\frac{d}{dt}{\Big|}_{t=0}R_t(z).$$
\end{enumerate}
 Following this terminology, we say that a curve $\gamma$ is a \emph{parabolic} (resp. \emph{hyperbolic}, \emph{rotational}) soliton of the ICSF if $\gamma$ is a soliton with respect to the vector field $X_1$ (resp. $X_2$ and $X_3$).

In the space $\mathfrak{X}(\h^2)$,  there is also a distinguished  vector field given by 
$$Y(x,y)=\partial_y.$$
This vector field is conformal, since its Lie derivative satisfies 
$$\mathcal{L}_{\partial_y}\langle,\rangle=-\frac{2}{y}\langle,\rangle.$$
 Observe that the integral curves of $Y$ are geodesics of $\h^2$, represented  by vertical lines. If we define $\Phi\colon\h^2\times\r\to\h^2$ by $\Phi(x,y,t)=(x,y+t)$, then $\frac{d}{dt}\Phi(x,y,t)=\partial_y$. 
Consider now a curve $\gamma(s)=(x(s),y(s))$ in $\h^2$.  We say that $\gamma$ is a soliton with respect to the vector field $Y$ if $\gamma_t(s)=\Phi(\gamma(s),t)=(x(s),y(s)+t)$ is a solution of \eqref{flow}.  Then, we have 
\begin{eqnarray*}
N(t,s)&=&(y(s)-t)\frac{(-y'(s),x'(s))}{\sqrt{x'(s)^2+y'(s)^2}}, \\
\kappa_t(s)&=&(y(s)-t)\kappa_e(s)+\frac{x'(s)}{\sqrt{x'(s)^2+y'(s)^2}}.
\end{eqnarray*} 
Since $ \frac{\partial \gamma_t}{\partial t}=(0,1)$, then 
$$\langle \frac{\partial \gamma_t}{\partial t},  N(t,s)\rangle = \frac{ x'(s)}{(y(s)-t)\sqrt{x'(s)^2+y'(s)^2}}.$$
Thus, \eqref{flow} is 
$$\frac{x'(s)}{y(s)+t} = -\frac{x'(s)^2+y'(s)^2}{x'(s)+(y(s)+t)\kappa_e\sqrt{x'(s)^2+y'(s)^2}}.$$
In particular, for $t=0$, we have that $\gamma$ satisfies
\begin{equation}\label{conformal}
\frac{1}{\kappa}=-\langle N,\partial_y\circ\gamma\rangle,
\end{equation}
In the theory of MCF for surfaces in hyperbolic space $\h^3$, solitons of   conformal vector fields  were studied in \cite{bl1,ma}. See also  \cite{bl2} when the ambient space is $\h^2\times\r$.

 
%%%%%%%%%%%%%%%%%%%%%%%
\section{Examples of the ICSF: case of constant curvature} \label{sec3}
%%%%%%
 
We now present several examples of the ICSF for curves of constant curvature. Note that the evolution of a geodesic under ICSF is not defined, since its curvature satisfies $\kappa=0$. Also, we say that the solution $\gamma = \gamma(t,s)$ of the ICSF is an \textit{ancient solution} if it is defined on the time interval of the form $\langle - \infty, T \rangle$, where $0 \leq T \leq \infty$. In other words, ancient solutions are defined for all times $t < 0$. 

\begin{example}[Flow of hyperbolic circles]
\label{ex:circle}
A hyperbolic circle $\gamma$ with center $(a,b) \in \r^2_+$ and radius $R$ coincides with the Euclidean circle whose center is located at $(a, \cosh(R))$ and whose radius is equal to $b\sinh(R)$. After applying an isometry of $\h^2$, we may assume, without loss of generality, that $(a,b)=(0,1)$. A parametrization of $\gamma$ is 
$$\gamma(s)=\sinh(R) \cdot (\cos(s),\sin(s)) + (0,\cosh(R)).$$
The curvature of $\gamma$ is $\kappa=\coth(R)$.

For the ICSF, it is natural to assume that the flow remains within the family of circles centered at $(0,1)$. Thus we consider the evolving family of curves
$$\gamma_t(s) = \sinh(r(t)) \cdot (\cos(s), \sin(s)) + (0,\cosh(r(t))).$$ 
We compute the terms appearing in the flow equation \eqref{soliton}. The curvature is given by $\kappa=\coth(r(t))$, while the unit normal vector $N$ and the time derivative $ \frac{\partial \gamma_t}{\partial t}$ satisfy
\begin{equation*}
\begin{split}
N &=- (\sinh(r)\sin(s) + \cosh(r)) \cdot (\cos(s),\sin(s)),\\
\frac{\partial \gamma_t}{\partial t} &= r' \cdot (\cosh(r)\cos(s),\sinh(r)+\cosh(r)\sin(s)).
\end{split}
\end{equation*}
Hence 
$$  \langle\left(\frac{\partial \gamma_t}{\partial t}\right)^\perp, N \rangle = -r'(t),$$
and so the equation \eqref{soliton} reduces to the following Cauchy problem:
\begin{equation*}
\begin{cases}
\displaystyle r' =  \frac{1}{\coth(r)}, \\
r(0) = R.
\end{cases}
\end{equation*}
The solution is $$r(t) = \sinh^{-1}(\sinh(R)e^{t}).$$
Since $r(t)\geq 0$ for all $t$, the ICSF exists for all $t \in \r $. In particular, this is an ancient solution of the ICSF. As $\lim_{t \to \infty} r(t) = \infty$, it follows that the hyperbolic circle expands infinitely under the ICSF. On the other hand, as $\lim_{t \to - \infty} r(t) = 0$ it follows that hyperbolic circles shrink to a point as $t \to -\infty$.
\end{example} 

\begin{example}[Flow of horocycles]
\label{ex:circle2}
Let $\gamma$ be the horocycle given by the Euclidean circle
$$\gamma(s)=R ( \cos(s),\sin(s))+(0,R),$$
which is tangent to the line $y=0$ at $(0,0)$. Assuming that the ICSF preserves the family of horocycles with ideal point $(0,0)$, we consider the evolving curves
 $$\gamma_t(s) = r(t) \cdot (\cos(s), \sin(s))+(0,r(t)).$$ 
Here $\kappa=1$, and 
$$\langle\left(\frac{\partial \gamma_t}{\partial t}\right)^\perp,  N \rangle = -\frac{r'(t)}{r(t)}.$$ 
Thus \eqref{soliton} becomes
\begin{equation*}
\begin{cases}
r' = r, \\
r(0) = R.
\end{cases}
\end{equation*}
The solution is $$r(t) = R\cdot e^t.$$
The ICSF is defined for all time $t\in\r$, so this is also an ancient solution of the ICSF. Since $\lim_{t\to\infty}r(t) = \infty$, we see that the horocycles also expand infinitely under the ICSF. 
\end{example} 


\begin{example}[Flow of equidistant lines]
\label{ex:circle3}
Let $0<c<R$. Consider the equidistant line given by the arc of the Euclidean circle 
$$\gamma(s)=R\left( \cos(s),\sin(s)\right)+(0,c),\qquad -\frac{c}{R}\leq \sin s\leq 1.$$
Its curvature is $\kappa=\frac{c}{R}$. We assume that the ICSF remains within the family of equidistant lines sharing the same ideal boundary point. Let $0<c(t)<r(t)$, and impose the condition 
$$r(t)^2-c(t)^2=R^2-c^2:=a^2,\quad a>0.$$
Define
$$\gamma_t(s) = r(t) \cdot (\cos(s),\sin(s))+(0, \sqrt{r(t)^2-a^2} ).$$ 
Then 
$$\kappa=\frac{\sqrt{r(t)^2-a^2}}{r(t)},\qquad \langle\left(\frac{\partial \gamma_t}{\partial t}\right)^\perp , N \rangle =- \frac{r'(t)}{\sqrt{r(t)^2-a^2}}.$$
Thus \eqref{soliton} reduces to the same ODE as before: 
\begin{equation*}
\begin{cases}
r' = r, \\
r(0) = R.
\end{cases}
\end{equation*}
Hence $$r(t) = Re^t.$$ 
The condition $r(t)^2>a^2$ requires the solution to satisfy
$$t > t_{\text{min}} := \log\left(\frac{a}{R}\right).$$

Since $t_{\text{min}} > - \infty$, this solution is not an ancient one. Also, note that at $t_{\text{min}}$ we have $r(t_{\text{min}}) = a$. Thus, as $t \to t_{\text{min}}$ the family $\gamma_t$ converges to the geodesic $s \mapsto a(\cos(s), \sin(s))$, which is precisely the geodesic to which the initial curve $\gamma$ is equidistant to. On the other hand, as $t \to \infty$ the equidistant line expands infinitely. 
\end{example} 

We summarize the previous examples in the following statement. 

\begin{proposition}
Let $\gamma$ be a curve in $\h^2$ with constant curvature $\kappa>0$.  
\begin{enumerate}
\item If $\gamma$ is a hyperbolic circle, then the ICSF is formed by hyperbolic circles sharing the same center as $\gamma$.   
\item If $\gamma$ is a horocycle, then the ICSF is formed by horocycles sharing the same ideal boundary point as $\gamma$. 
\item If $\gamma$ is an equidistant line, then the ICSF is formed by equidistant lines sharing the same ideal boundary point as $\gamma$. 
\end{enumerate}
Furthermore, the ICSF of all of the curves above is expanding, while the backward time limits are the common center, the ideal boundary point and a geodesic curve respectively.
\end{proposition}

  
 %%%%%%%
 \section{Parabolic solitons of ICSF}\label{sec5}
%%%%%%%%%%%%%%%%%%%%%%%%%%

In this section, we study parabolic solitons of the ICSF. We point out that two-dimensional parabolic solitons of the MCF  were already classified in \cite{bl1}. 

Let $\gamma(s)=(x(s),y(s))$ be a curve in $\h^2$, parametrized by Euclidean arc length. Then there is a smooth function $\theta=\theta(s)$ such that 
$$x'(s)=\cos\theta(s),\qquad y'(s)=\sin\theta(s).$$
Notice that $\kappa_e=\theta'$ and from \eqref{ke2}, we have 
$$\kappa=y\kappa_e+\cos\theta.$$
Moreover, $N=y(-\sin\theta,\cos\theta)$, so 
$$\langle N,X_1\rangle=-\frac{\sin\theta}{y}.$$
Equation \eqref{soliton} is therefore equivalent to

\begin{equation}
\label{para1}
\begin{cases}
x' = \cos\theta \\
y' = \sin\theta \\
\displaystyle \theta' =  \frac{1}{\sin\theta} - \frac{\cos\theta}{y}.
\end{cases}
\end{equation}

A consequence of the fact that $\sin\theta(s)$ cannot take the value $0$ is the following.

\begin{proposition}\label{prg}
 Any parabolic soliton to the ICSF is a graph on the $y$-axis.
\end{proposition}

Our first result establishes the fact that parabolic solitons of the ICSF cannot have constant curvature.

\begin{proposition}
There are no parabolic solitons of the ICSF with constant curvature.
\end{proposition}

\begin{proof} 
Let $\gamma(s)=(x(s),y(s))$ be a parabolic soliton with constant curvature. Then $\gamma$ defines $\theta(s)$ such that $(x(s),y(s),\theta(s)($ is a  solution of system \eqref{para1}. Then the Euclidean curvature $\kappa_e = \theta'$ is constant. We distinguish between two cases:
\begin{enumerate}
\item Case $\theta' = 0$. Then $\theta(s) = \theta_0$ for some constant $\theta_0\in\r$, and the soliton $\gamma$ can be written as 
 $$\gamma(s) = s  (\cos\theta_0 , \sin\theta_0 ) + (a,b),$$
 with $a,b\in\r$ and  $\sin\theta_0\not=0$. 
Substituting into the last equation of \eqref{para1}, we obtain 
$$ s \sin\theta_0 +b = \sin\theta_0   \cos\theta_0,$$
which implies $\sin\theta_0 = 0$, but this is a contradiction. Therefore, this case is impossible.

\item Case $\theta' = c\not=0$. Then $\theta(s) = cs$ and the soliton can be written as 
$$\gamma(s) = \frac{1}{c} (\sin(cs), -\cos(cs)) + (x_0,y_0).$$
Substituting into \eqref{para1} yields
$$ c y_0 \sin(cs) + \frac{\cos(cs)}{c} - y_0=0.$$
Since this holds for all $s\in\r$, we deduce $1/c=0$, which is a contradiction. Therefore, this case is impossible as well. 
\end{enumerate}
\end{proof}


Now we want to study geometric properties of parabolic solitons. Let $\gamma(s)=(x(s),y(s))$ be a solution of \eqref{para1}. Note that the function $x$ is completely determined by $\theta$ so it suffices to solve \eqref{para1} for functions $y,\theta$. Therefore, the system \eqref{para1} is equivalent to the system

\begin{equation}
\label{para2}
\begin{cases}
y' = \sin\theta \\
\displaystyle \theta' =  \frac{1}{\sin\theta} - \frac{\cos\theta}{y},
\end{cases}
\end{equation}


Since the system \eqref{para2} is   independent of $x$,   we can prescribe the initial condition $x(0)$ without loss of generality. Moreover, there are no equilibrium points.

In order to determine the phase space associated to \eqref{para2}, we first note the following result, which is a standard consequence of the uniqueness theorem for ODEs,  and which tells us that we can consider $\theta$ to be contained within the interval $(0,\pi)$.

\begin{proposition}\label{pr52}
If $\alpha(s) = (y(s),\theta(s))$ is a solution of \eqref{para2}, then  $$\overline{\alpha}(s) = (y(-s),\theta(-s) - \pi)$$ is also a solution of \eqref{para2}. 
\end{proposition}

Therefore, the phase space associated to the system \eqref{para2} is equal to the set 
$$\Theta = \left \{ (y, \theta) \colon y > 0, \theta \in (0,\pi) \right\}.$$
To determine geometric properties of parabolic solitons, we are going to be studying the orbits of the system \eqref{para2}. Let $\alpha(s) = (y(s), \theta(s))$ be any solution of the system \eqref{para2} defined on the interval $I = (s_m, s_M)$, which we assume to be maximal.

Each solution $\gamma(s)$ of \eqref{para1}, with its corresponding function $\theta(s)$,  has an associated unique orbit $\alpha_\gamma(s)= (y(s),\theta(s))$ of \eqref{para2}. Standard ODE theory ensures that through every point of $\Theta$ passes a unique orbit, and distinct orbits do not intersect. 
Let us observe that the system \eqref{para2} has no equilibrium point in $\Theta$ and thus, a phase space analysis of the dynamics of the orbits using these points cannot be done.






The points in which the Euclidean curvature vanishes are determined by $\theta' = 0$, and they all lie on the curve 
$$\Gamma\colon y = \frac{1}{2} \sin(2\theta) \subseteq \Theta.$$
 Since $y>0$, the curve $\Gamma$ appears in the phase space $\Theta$ only when $\theta \in (0, \pi/2)$.

Using the curve $\Gamma$, we can determine the regions of the phase space in which the coordinate functions $y,\theta$ of orbit $\alpha$ are strictly monotone. First of all, since $\theta \in (0, \pi)$ we have that $y' = \sin \theta > 0$ so $y$ is strictly increasing on the whole phase space. As for the function $\theta$, we have that
$$\theta' = \frac{1}{\sin\theta} - \frac{\cos\theta}{y} = \frac{y-\frac{1}{2}\sin(2\theta)}{y\cdot \sin\theta}.$$
Thus, the monotonicity regions are

\begin{equation*}
\left\{\begin{split}
& \theta \in (0, \frac{\pi}{2} ) \colon  \left\{\begin{array}{ll}
\theta \text{ increases if and only if} &  y >\frac{1}{2}\sin(2 \theta), \\
\theta \text{ decreases if and only if} &  y < \frac{1}{2}\sin(2 \theta).\\
\end{array}\right.\\
& \theta \in ( \frac{\pi}{2},\pi) \colon\mbox{$\theta$ is increasing}.
\end{split}\right.
\end{equation*}
 
 In the study of solutions of \eqref{para2}, we fix initial conditions. As we have seen,   different choices for  $x(0)$ give the same solution $\gamma(s)=(x(s),y(s))$, up to a horizontal translation. Thus, we can assume, without loss of generality, that $x(0)=0$. We fix the initial condition on $\theta $ to be $\theta(0)=\pi/2$.  Later, we will analyze other initial conditions for the function $\theta = \theta(s)$. The classification is obtained in Theorems \ref{t55} and \ref{t57}, respectively. In both results, the geometric properties of the parabolic solitons  depend on the initial condition $y(0)=y_0$. 
 
 
 
%%%%%% 
\subsection{Solutions of \eqref{para1} with initial condition $\theta(0) = \pi / 2$}

Let $\gamma(s)=(x(s),y(s))$ be a solution of \eqref{para1} with initial conditions 
$$x(0)=0,\quad y(0)=y_0,\quad \theta(0)=\frac{\pi}{2}.$$
Let $\alpha_\gamma(s) = (y(s), \theta(s))$ be the corresponding orbit of the system \eqref{para2}, with the initial condition $\alpha_\gamma(0) = (y_0, \pi / 2)$. 


\begin{lemma}\label{le54}
    Let $\alpha = (y,\theta)$ be an orbit of the system \eqref{para2} with the initial condition $\alpha(0) = (y_0, \pi/2)$. Then $\lim_{s \to s_M} \alpha(s) = (Y, \pi)$, where $Y < \infty$.
\end{lemma}

\begin{proof}
At $s = 0$, we have that $\theta'(0) = 1 > 0$ so there is an interval around $s = 0$ on which the function $\theta$ is increasing ant, thus, for $s>0$ close to $0$,   the orbit $\alpha_\gamma(s)$ is contained in  $\theta \in (\pi/2, \pi)$. This implies that  the functions $y$ and $\theta$ are increasing for $s > 0$. Let $\lim_{s\to s_M}(y(s),\theta(s))=(Y,\theta^\ast)$. 



    By Proposition  \ref{prg}, we know that   $\theta$ is a function of $y$, i.e. $\theta = \theta(y)$. Therefore, we have
    $$\frac{d \theta}{dy} = \frac{d \theta}{ds}   \left( \frac{dy}{ds} \right)^{-1} = \frac{1}{\sin\theta}  \left( \frac{1}{\sin\theta} - \frac{\cos\theta}{y} \right) = \frac{1}{\sin^2 \theta} - \frac{\cot\theta}{y}.$$
    Note that $d\theta / dy$ and $d\theta / ds$ have the same sign, so the regions where those functions are monotone coincide.

    Fix some $\overline{y} > y_0$ and denote $\overline{\theta} = \theta(\overline{y})$. Since $d\theta / dy > 0$ on $(y_0, Y)$, we have that $\theta(y)$ is an increasing function on that interval. Because of that we have $\theta(y) \in (\pi/2, \pi)$ for all $y > \overline{y}$ and $\theta(y) > \overline{\theta}$ for all $y > \overline{y}$. Since the function $\theta \mapsto \cot\theta$ is negative and decreasing on $(\pi/2, \pi)$, for all $y > \overline{y}$ we have  $ \cot\overline{\theta} > \cot \theta$. If $m:= - \cot\overline{\theta}>0$, then 
    $$ -\frac{\cot\theta}{y} > \frac{m}{y}.$$
  This yields
    $$\frac{d\theta}{dy} = \frac{1}{\sin^2 \theta} - \frac{\cot\theta}{y} > \frac{m}{y}$$
    for all $ y > \overline{y}$.   Integrating   from $\overline{y}$ to $y$, $y > \overline{y}$, we have $$\theta(y) - \overline{\theta} \geq m\log(y) - m\log(\overline{y}).$$
    If $Y$ were $\infty$, letting $y \to \infty$ we get $\pi \geq \lim_{y \to \infty} \theta(y)\geq \infty$, which is a contradiction. This proves that $Y < \infty$.

    On the other hand, if it were $\lim_{y \to Y} \theta(y) = \theta^\ast < \pi$, then also $\lim_{s \to s_M} \theta(s) = \theta^\ast$. Since $Y < \infty$, we could extend the orbit past $Y$ by considering the initial condition to be equal to $(Y, \theta^\ast)$, but this is impossible since the interval $I$ is maximal. Therefore, $\lim_{s \to s_M} \theta(s) = \lim_{y \to Y} \theta(y) = \pi$.

\end{proof}

 On the other hand, for  $s < 0$ and near $0$, the orbit $\alpha_\gamma(s)$ is contained in the region $\theta \in (0, \pi / 2)$ so there are two possible cases:
\begin{enumerate}
    \item If $\alpha_\gamma$ intersects the curve $\Gamma$ at some $\overline{s} < 0$, then $\theta(s)$ is decreasing on the interval $(s_m, \overline{s})$, here $\lim_{s \to s_m} \alpha_\gamma(s) = (0, \theta_\ast)$.
    \item If $\alpha_\gamma$ does not intersect the curve $\Gamma$, then $\theta(s)$ is increasing on the whole interval $(s_m, 0)$ and so $\lim_{s \to s_m} \alpha_\gamma(s) = (\bar{Y},0)$, with $\bar{Y}>0$.
\end{enumerate}



\begin{figure}[h!t]
\centering
\includegraphics[width=0.6\linewidth]{phase-parabolic.pdf}
\caption{Orbits of the system \eqref{para2} varying values of $y(0)$. The black line is the curve $\Gamma$.}
\label{fig:parabolic-orbit-pi-over-2}
\end{figure}
If the orbit $\alpha_\gamma$ intersects (resp. does not intersect) the curve $\Gamma$, then we will refer to that orbit, and the corresponding soliton, as a \textit{Type I} (resp. \textit{Type II})  orbit/parabolic soliton.  

 








The following lemma tells us that if the orbit $\alpha_\gamma$ intersects $\Gamma$ for some $\hat{s}$, then it is contained inside the region determined by $\Gamma$ for all $s \in (s_m, \hat{s})$. 

\begin{lemma}
\label{lem-unique-intersection}
    Let $\alpha_\gamma = (y, \theta)$ be an orbit of the system \eqref{para2} with the initial condition $\alpha_\gamma(0) = (y_0, \pi / 2)$. If $\alpha_\gamma$ is of Type I, then there exists unique $\hat{s} < 0$ such that $\alpha_\gamma(\hat{s}) \in \Gamma$.
\end{lemma}

\begin{proof}
    We know that $\alpha_\gamma$ meets $\Gamma$  at least one point. Without loss of generality, let   that $\hat{s}$ is the largest value of parameter $s$ such that $\alpha_\gamma(s)$ intersects $\Gamma$. By contradiction, suppose that there is some $\overline{s} < \hat{s}$ such that $\alpha_\gamma(\overline{s}) \in \Gamma$ as well. Denote $\overline{y} = y(\overline{s})$ and $\hat{y} = y(\hat{s})$. Note that $\overline{y} < \hat{y}$ since $y$ is an increasing function, and note that $\alpha_\gamma(s)$ is contained inside the region determined by $\Gamma$ for all $s \in (\overline{s}, \hat{s})$. Because of that, $\theta(y)$ is decreasing on the interval $(\overline{y}, \hat{y})$.

    Since $\frac{d\theta}{dy}(\overline{y}) = 0$, there are two possible cases:

    \begin{enumerate}
	\item  $\overline{y}$ is a local extrema of the function $\theta(y)$. Since $\theta(y)$ is decreasing on the interval $(\overline{y}, \hat{y})$, the point $\overline{y}$ has to be a local maximum. Hence, there exists an interval on which $\alpha_\gamma(s)$ is contained inside the region determined by $\Gamma$ and on which the function $\theta$ is increasing, but this is impossible since $\theta$ is decreasing whenever $\alpha_\gamma$ is inside $\Gamma$. 
	
	\item $\overline{y}$ is an inflection point of the function $\theta(y)$. Hence, for $y < \overline{y}$, the orbit $\alpha_\gamma(y)$ should be outside $\Gamma$ and the function $\theta(y)$ should be decreasing. This is impossible since $\theta$ is increasing in that case.   
    \end{enumerate}
\end{proof}

\begin{lemma}\label{le56}
    Let $\alpha_\gamma = (y, \theta)$ be an orbit of the system \eqref{para2} with the initial condition $\alpha_\gamma(0) = (y_0, \pi/2)$. If $\alpha_\gamma$ is of Type I, then $\lim_{s \to s_m} \alpha_\gamma(s) = (0, \pi/2)$.
\end{lemma}

\begin{proof}
    Since $\alpha_\gamma$ is a Type I orbit, and by Lemma {lem-unique-intersection}, there exists a unique $\hat{s} < 0$ such that $\alpha_\gamma(\hat{s}) \in \Gamma$. Denote $\hat{y} := y(\hat{s})$ and $\hat{\theta} := \theta(\hat{s})$. Moreover, we know that  $\alpha_\gamma$ is contained inside the region determined by $\Gamma$ for all $y \in (0, \hat{y})$. Then $\lim_{s \to s_m} \alpha_\gamma(s) = (0, \theta_\ast)$, where $\theta_\ast \in (\hat{\theta}, \pi / 2] \subseteq (0, \pi/2]$.

    Suppose that $\theta_\ast \neq \pi / 2$. By Proposition \ref{prg}, we  can consider $\theta$ as a function of $y$, i.e. $\theta = \theta(y)$ for $y \in (0, \hat{y})$. Since $\alpha_\gamma$ is contained inside $\Gamma$, we know that $\theta(y)$ is a decreasing function on the interval $(0, \hat{y})$. Therefore, we have $\theta(y) \in (\hat{\theta}, \theta_\ast)$ for all $y \in (0,\hat{y})$.

    Since the functions $\theta \mapsto \frac{1}{\sin^2 \theta}$ and $\theta \mapsto \cot\theta$ are positive and continuous on $(0, \pi/2)$, they are in particular bounded on $(0,\hat{y})$, since for $y \in (0, \hat{y})$, we have $\theta(y) \in (\hat{\theta}, \theta_\ast)$ and $\theta_\ast \neq \pi / 2$. So, there are constants $A,B > 0$ such that
    \begin{equation*}
       	\begin{split}
	    \frac{1}{\sin^2 \theta(y)} &\leq A, \quad  y \in (0, \hat{y}) \\
	    \cot(\theta(y)) &\geq B, \quad  y \in (0, \hat{y}).
       	\end{split}
    \end{equation*}

    In particular, for all $y \in (0, \hat{y})$ we have
    $$\frac{d\theta}{dy} = \frac{1}{\sin^2\theta} - \frac{\cot\theta}{y} \leq A - \frac{B}{y}.$$

    Integrating from $y$ to $\hat{y}$, $y < \hat{y}$, we get
    $$\theta(y) - \hat{\theta} \leq A(\hat{y} - y) -B\log(\hat{y}) + B\log(y).$$

    Letting $y \to 0,$ we get $$\theta_\ast - \hat{\theta} \leq - \infty,$$ which is an obvious contradiction. Hence, $\theta_\ast = \pi/2$ and so $$\lim_{s \to s_m} \theta(s) =\frac{ \pi}{2}.$$
\end{proof}


We are in conditions to describe geometrically the parabolic solitons of the ICSF. 

\begin{theorem}\label{t55}
Let $\gamma$ be a parabolic soliton of the ICSF solving \eqref{para1} with initial conditions $(x(0),y(0),\theta(0))=(0,y_0,\pi/2)$. Then $\gamma$ is a graph on the $y$-axis. Moreover, there exists a constant $H > 0$ such that:
\begin{enumerate}
\item  If $0<y_0<H$, then $\gamma$ approaches the line $y=0$ orthogonally, and there exists a constant $Y > 0$ such that the line $y = Y$ is a horizontal asymptote of $\gamma$ (type I);
\item If $y_0>H$, then $\gamma$ is a graph on a bounded interval $(m, M)$ on the $y$-axis (type II). Writing $\gamma$ as $x=x(y)$, then   $x(y)$ is concave and satisfies 
$$\lim_{y\to M}x'(y)=-\infty,\quad \lim_{y\to m}x'(y)=\infty.$$
\end{enumerate}
\end{theorem}

\begin{proof}
By Proposition \ref{prg}, we know that $\gamma$ is a graph on the $y$-axis. Lemma \ref{le54} implies that one end of $\gamma$ is asymptotic to the horizontal line $y=Y$.
Let $\alpha_\gamma$ be the  orbit  corresponding to $\gamma$. Since an orbit may or may not intersect $\Gamma$, there is $H> 0$ such that $\alpha_\gamma$ intersects $\Gamma$ whenever $y_0 < H$, while for $y_0 > H$  the orbit  $\alpha_\gamma$ and $\Gamma$ are disjoint. 

We now distinguish two cases:
\begin{enumerate}
\item Case   $y_0 < H$. Then  $\alpha_\gamma$ is of Type I. By Lemma \ref{le56}, the limit of $\alpha_\gamma$ as $s\to s_m$ implies that one end of $\gamma$ goes to $y=0$ and that intersection with this axis is orthogonal. 



\item Case $y_0>H$. Then $\alpha_\gamma$ is of Type II. We have seen that in such a case,  $\lim_{s \to s_m} \alpha_\gamma(s) = (\bar{Y},0)$, with $\bar{Y}>0$. This implies that one end of $\gamma$ is asymptotic to the line $y=\bar{Y}$. In particular, $\gamma$ is a graph on a  bounded interval $(\bar{Y},Y)$. 

 
To prove that $x=x(y)$ is concave, by the chain rule and \eqref{para1}, we have 
$$x''(y)=-\frac{\theta'}{\sin^3\theta}=-\frac{y-\sin\theta\cos\theta}{y\sin^4\theta}<0$$
because $\alpha_\gamma\cap\Gamma=\emptyset$, and thus, $y>\sin\theta(s)\cos\theta$. This proves that $\gamma$ is a concave graph on the $y$-axis.  Finally, we have 
$$\lim_{y \to Y} x'(y) =\lim_{s\to s_M}\frac{\cos\theta(s)}{\sin\theta(s)}=-\infty,\quad \lim_{y \to \bar{Y}} x'(y) =\lim_{s\to s_m}\frac{\cos\theta(s)}{\sin\theta(s)}=\infty .$$
\end{enumerate}
\end{proof}

Examples of the  types I and II of parabolic solitons to the ICSF are shown in Figure  \ref{fig:parabolic-solitons}.


%%%%%%%%
\subsection{Solutions of \eqref{para1} with initial condition $\theta(0) \neq \pi / 2$}

We describe the orbits $\alpha_\gamma$ with the initial condition $\theta(0) \neq \pi/2$. 

\begin{lemma}\label{le58}
    Let $\alpha_\gamma = (y, \theta)$ be an orbit of the system \eqref{para2} with the initial condition $\theta(0) < \pi /2$. Then $\alpha_\gamma$ can be considered as either Type I or Type II orbit.
\end{lemma}

\begin{proof}
    We distinguish between three cases: 
    \begin{enumerate}
	\item If $y(0) < \frac{1}{2}\sin(2\theta(0))$, then $\alpha_\gamma(0)$ is contained inside the region determined by the curve $\Gamma$. Using the same arguments as before, we conclude that $$\lim_{s \to s_m} \alpha_\gamma(s) = (0, \pi/2), \quad \lim_{s \to s_M} \alpha_\gamma(s) = (Y, \pi),$$
		    i.e. $\alpha_\gamma$ is a Type I orbit.
	\item If $y(0) = \sin(2\theta(0))$, then $\alpha_\gamma(0) \in \Gamma$. Since the system \ref{para2} has no equilibrium points, $\theta$ is not a constant function. Hence, $y(0)$ is a critical point of $\theta(y)$, and so $\alpha_\gamma$ is a Type I orbit.

	\item If $y(0) > \sin(2\theta(0)$, then $\alpha_\gamma(0)$ is contained in the outer region determined by the curve $\Gamma$. Because of that, $\theta(y)$ is strictly increasing. We now claim that there exists $\overline{y} > y_0$ such that $\theta(\overline{y}) = \pi / 2$.

	    Assume that $\theta(y) < \pi/2$ for all $y > y_0$. Since $\theta(y)$ is an increasing function, that is only possible if the line $\theta = \pi/2$ is a horizontal asympotote of the function $\theta(y)$, i.e. if
	    $$Y = \infty, \quad \lim_{y \to \infty} \theta(y) = \pi/2, \quad \lim_{y \to \infty} \frac{d \theta}{dy} = 0.$$

	    In that case, we have the following:
$$      	    0 = \lim_{y \to \infty} \frac{d \theta}{dy} = \lim_{y \to \infty} \frac{1}{\sin^2 \theta} - \frac{\cot\theta}{y}   = 1,$$
	    but this is an obvious contradiction. Hence, there exists some $\overline{y} > y_0$ such that $\theta(\overline{y}) = \pi / 2$, so $\alpha_\gamma$ can be considered as an orbit with the initial condition $\theta(0) = \pi/2$. In particular, $\alpha_\gamma$ is either a Type I or Type II orbit.
    \end{enumerate}
\end{proof}


\begin{lemma}\label{le59}
    Let $\alpha_\gamma = (y, \theta)$ be an orbit of the system \eqref{para2} with the initial condition $\theta(0) > \pi/2$. Then $\alpha_\gamma$ satisfies one of the following:
    \begin{itemize}
	\item[i)] $\alpha_\gamma$ is a Type I or a Type II orbit
	\item[ii)] $\theta(y) > \pi/2$ for all $y > 0$ and $$\lim_{s \to s_m} \alpha_\gamma(s) = (0, \pi/2), \quad \lim_{s \to s_M} \alpha_\gamma(s) = (Y, \pi)$$.
    \end{itemize}
\end{lemma}

\begin{proof}
   	If the orbit $\alpha_\gamma(s)$ intersects the line $\theta = \pi/2$, then $\alpha_\gamma$ can be considered as an orbit with initial condition $\theta(0) = \pi/2$, so $\alpha_\gamma$ is an orbit of Type I or Type II.

	So, let us assume that $\alpha_\gamma$ does not intersect the line $\theta = \pi/2$. Since $\theta(y)$ is an increasing function, it follows that $\lim_{s \to s_M} \alpha_\gamma(s) = (Y, \pi)$.

	On the other hand, we claim that $\lim_{y \to 0} \theta(y) = \pi/2$. In order to prove our claim, suppose that $\lim_{y \to 0} \theta(y) = \theta_\ast > \pi/2$. Fix some $\overline{y} > 0$ and denote $\overline{\theta} = \theta(\overline{y})$. Since $\theta_\ast > \pi/2$, the function $\theta \mapsto \cot\theta$ is negative and continuous on the interval $(0, \overline{y})$ because for all such $y$ we have that $\theta(y) \in (\theta_\ast, \overline{\theta}) \subseteq (\pi/2, \pi)$. So there exists some $m < 0$ such that for all $y \in (0, \overline{y})$ we have
	$$  - \frac{\cot\theta}{y} \geq - \frac{m}{y}.$$

	Because of that, we have the bound
	$$\frac{d\theta}{dy} = \frac{1}{\sin^2(\theta)} - \frac{\cot \theta}{y} \geq - \frac{m}{y}.$$

	Integrating from $y$ to $\overline{y}$, $y < \overline{y}$, we get
	$$\overline{\theta} - \theta(y) \geq -m\log(\overline{y}) + m\log(y).$$

	Letting $y \to 0$, it follows that $\overline{\theta} - \theta_\ast \geq \infty$, which is an obvious contradiction. Hence, $\lim_{y \to 0} \theta(y) = \pi/2$.
\end{proof}

Using the previous two lemmas, we have the geometric description of the parabolic solitons when  $\theta(0)\not=\frac{\pi}{2}$.  


  \begin{theorem}\label{t57}
  Let $\gamma$ be a parabolic soliton of the ICSF solving \eqref{para1} with initial conditions $(x(0),y(0),\theta(0))=(0,y_0,\theta(0))$,   with $\theta(0)\not=\frac{\pi}{2}$. Then $\gamma$ belongs to one of the two types described in  Theorem \ref{t55} or $\gamma$ is a concave graph defined at some bounded domain $(0,x_2)$, where the intersection of the branch of $\gamma$ which meets $x$-axis is orthogonal.  
 \end{theorem}
 
 \begin{proof} By Lemmas \ref{le58} and \ref{le59}, the parabolic soliton can be of Types I or II of Theorem \ref{t55}. Otherwise, Lemma \ref{le59} implies that $\theta(s)>\pi/2$ for all $y>0$. In such a case, by the properties of $\alpha_\gamma$, we know that $\gamma$ arrives to the $x$-axis and that this intersection is orthogonal. Concavity of $\gamma$ is proved as in Theorem \ref{t55} because $y>\sin\theta\cos\theta$ since $\theta\in (\frac{\pi}{2},\pi)$. 
  
 \end{proof}

The last case in Theorem \ref{t57} is shown in Figure   \ref{fig:parabolic-solitons}, right.




 




\begin{figure}[hbtt]
 \begin{center}
 \includegraphics[width=.13\textwidth]{parabolic11.pdf}\qquad\qquad \includegraphics[width=.22\textwidth]{parabolic12.pdf}\qquad \qquad\includegraphics[width=.25\textwidth]{parabolic13.pdf}
\end{center}
\caption{Parabolic solitons of the ICSF. Left: $\theta(0)=\frac{\pi}{2}$, $y(0)=1$. Middle: $\theta(0)=\frac{\pi}{2}$, $y(0)=2$. Right: $\theta(0)=\frac{3\pi}{4}$, $y(0)=0.2$. }
\label{fig:parabolic-solitons}
\end{figure}


%%%%%%%%%%%%%%%%%%%%%%%%%%%%%%%%%%%%%%%%%%%%%%%%%%%%%%%%%%%%%%%
\section{Conformal solitons of the ICSF}\label{sec6}
%%%%%%%%%%%%%%%%%%%%%%%%%%%%%%%%%%%%%%%%%%%%%%%%%%%%%%

In this section, we study conformal solitons for the ICSF in $\h^2$. Let $\gamma(s) = (x(s), y(s))$ be a curve in $\h^2$, parametrized by Euclidean arc-length. If $\gamma$ is a conformal soliton of the ICSF, then it is characteried by \eqref{conformal}. Using the same approach as in the previous section, and if $\theta$ is the angle function of $\gamma$, equation \eqref{conformal} is equivalent to the following system
\begin{equation}\label{system-conformal}
\left\{\begin{split}
x'&=\cos\theta\\
y'&=\sin\theta\\
\theta'&=-\frac{1}{\cos\theta}-\frac{\cos\theta}{y}.
\end{split}
\right.
\end{equation}

Firstly, we prove that conformal solitons cannot have constant curvature.

\begin{proposition}
   There are no conformal solitons of the ICSF in $\h^2$ with constant curvature.
\end{proposition}

\begin{proof}
   Since the hyperbolic curvature is constant if and only the Euclidean curvature is constant, then the function $\theta'$ is constant.  
   \begin{enumerate}
       \item Case $\theta' = 0$. Then $\gamma$ can be  parametrized as 
       $$\gamma(s) = s   (\cos\theta_0, \sin\theta_0) + (a,b).$$ 
	   The third equation of the system \eqref{system-conformal} is then equivalent to $s   \sin\theta_0 + \cos^2\theta_0 + b=0$. Thus,  	         $\sin\theta_0 = 0$ and $ \cos^2\theta_0 + b = 0$. This gives $b=-1$. However, we have  $\gamma(s)=(s\cos\theta_0+a,b)$ and thus, $b$ should be positive. This is a contradiction.  

       \item Case $\theta' = c \neq 0$. Then $\gamma$ can be  parametrized as $$\gamma(s) = \frac{1}{c} (\sin(cs), -\cos(cs)) + (a, b).$$
	   The third equation of the system \eqref{system-conformal} is then equivalent to  $ \cos(cs)    ( \frac{1}{c} - cb  ) - b=0$. Thus $b=0$ and $\frac{1}{c}-cb=0$. This is a contradiction.  
   \end{enumerate}
\end{proof}

We   study the geometric properties of conformal solitons  using the same approach as in Section \ref{sec5}. A first result is   the following result on the symmetry of conformal solitons, which is analogous to Proposition \ref{pr52}.

\begin{proposition} 
\label{prop:sym}
Let $(x(s), y(s), \theta(s))$ be a solution of the system \eqref{system-conformal} with initial conditions $(x_0, y_0, \theta_0)$. Then   $(\bar{x}(s),\bar{y}(s),\bar{\theta}(s))=(2x_0-x(-s), y(-s), -\theta(-s))$ is also a solution of  \eqref{system-conformal} with initial conditions $(x_0,y_0,-\theta_0)$. 
\end{proposition}
 
 Note that the function $x(s)$ does not appear in the third equation of \eqref{system-conformal}. Thus, and without loss of generality, we can take $x(0)=0$ as the initial condition. Again, we  study the solutions of the   system
\begin{equation}
\label{system-conformal-reduced}
\begin{cases}
   y' = \sin\theta \\
   \displaystyle \theta ' = - \frac{1}{\cos\theta} - \frac{\cos\theta}{y}
\end{cases}
\end{equation}

In the remainder of this section, if $\gamma$ is a conformal soliton then we will denote by $\alpha_\gamma$ the orbit of the system \eqref{system-conformal-reduced} induced by $\gamma$. The next proposition gives us some properties of the orbits $\alpha_\gamma$, and it is a consequence of uniqueness of ODEs.

\begin{proposition}
    Let $\alpha_\gamma(s)=(y(s),\theta(s))$ be an orbit of the system \eqref{system-conformal-reduced}.
    \begin{enumerate}
	\item $\overline{\alpha_\gamma}(s) = (y(-s), - \theta(-s)$ is also an orbit of the system \eqref{system-conformal-reduced}. Consequently, $\alpha_\gamma$ is symmetric with respect to the line $\theta = 0$.

	\item if $\theta \in (-\pi/2, \pi / 2)$, then $\overline{\alpha_\gamma}(s) = (y(s), -\theta(s) + \pi)$ is an orbit of the system \eqref{system-conformal-reduced} such that $-\theta(s) + \pi \in (\pi / 2, 3\pi/2)$.
       
    \end{enumerate}
\end{proposition}

Therefore, the phase plane of   \eqref{system-conformal-reduced} is  
$$\Theta= \left\{(y,\theta)\colon y>0,\theta\in(-\frac{\pi}{2},\frac{\pi}{2}) \right \}.$$
It is immediate that there are no equilibrium points. Since $\cos\theta(s)\not=0$, then we have the analogous of Proposition \ref{prg}.

\begin{proposition}\label{prg2}
 Any conformal soliton to the ICSF is a graph on the $x$-axis.
\end{proposition}
 

\begin{lemma}\label{le64}
    Let $\alpha_\gamma = (y,\theta)$ be an orbit of the system \eqref{system-conformal-reduced}. Then $\theta'$ has the same sign on the whole interval $I$.
\end{lemma}

\begin{proof}
If at some point  $s \in I$, we have  $\theta'(s) = 0$, then   the last equation of  \eqref{system-conformal-reduced} gives $ y(s) = - \cos^2(\theta(s)) < 0$, which is not possible.  
\end{proof}

 

In the following two subsections, we analyze the   orbits (and hence the corresponding solitons) of \eqref{system-conformal-reduced} depending on the initial condition $\theta(0)$.  In Fig.  \ref{fig:conformal-orbit-0}, we show orbits of the system \eqref{system-conformal-reduced}.

%%%%%%%%%%%
\subsection{Solutions of \eqref{system-conformal-reduced} with initial condition $\theta(0) = 0$.}

Let $\alpha_\gamma = (y, \theta)$ be an orbit of the system \eqref{system-conformal-reduced} with initial conditions $(y(0),\theta(0)) = (y_0,0)$.   Denote by $(s_m,s_M)$ the maximal domain of $\alpha_\gamma$. Since $\theta'(0) = -1 - \frac{1}{y_0}$ and from Lemma \ref{le64},   it follows that $\theta$ is a strictly decreasing function. Since  $\theta \in (0, \pi / 2)$,  we have that $y' = \sin\theta > 0$, hence $y(s)$ is an increasing function. Therefore, for  $s < 0$ and near to $0$, the orbit $\alpha_\gamma$ lies contained in the region $\theta \in (0, \pi/2)$.  so $\lim_{s \to s_m} \alpha_\gamma(s) = (y_\ast, \theta_\ast)$, where $\theta_\ast \in (0, \pi/2 ]$. Because of the symmetry, we get $\lim_{s \to s_M} \alpha_\gamma(s) = (y_\ast, - \theta_\ast)$.



\begin{figure}[h]
\centering
\includegraphics[width=0.6\linewidth]{phase-conformal2.pdf}
\caption{Orbits of the system \eqref{system-conformal-reduced}   varying values of  $y(0)$.}
\label{fig:conformal-orbit-0}
\end{figure}


\begin{lemma}
\label{lem-conformal-angles} We claim  $\theta_\ast = \pi/2$.
\end{lemma}

\begin{proof}
   Suppose that $\theta_\ast \in (0, \pi/2)$. Then it must be that $y_\ast = 0$, since otherwise we could extend the solution past $s_m$. Since $y' = \sin\theta \neq 0$, we can write $\theta$ as a function of $y$, i.e. $\theta = \theta(y)$. Then
   $$\frac{d\theta}{dy} = \frac{d\theta}{ds}   \left( \frac{dy}{dy}\right)^{-1} = \frac{1}{\sin\theta}  \left(- \frac{1}{\cos\theta} - \frac{\cos\theta}{y} \right) < 0,$$
   so $\theta = \theta(y)$ is a decreasing function.

   Fix some $\hat{y} < y_0$ and denote $\hat{\theta} = \theta(\hat{y}) > 0$. Also, note that $$\frac{d\theta}{dy} = - \frac{2}{\sin(2\theta)} - \frac{\cot(\theta)}{y}.$$

   Since $\theta_\ast < \pi/2$, we have that $2\theta(y) \in (2\hat{\theta}, 2\theta_\ast)$ for $y \in (0, \hat{y})$. So the function $\theta \mapsto - \frac{2}{\sin(2\theta)}$ is negative on $(0,y_0)$. Also, since $\theta(y) \in (\hat{\theta}, \theta_\ast)$ for all $y \in (0, \hat{y})$, the function $\theta \mapsto -\cot\theta$ is negative and bounded on the interval $(0, \hat{y})$, so there exists some $m < 0$ such that $$-\cot\theta < m, \quad y \in (0, \hat{y}).$$

   Taking these bounds into account, it follows that on the interval $(0, \hat{y})$ we have $$\frac{d\theta}{dy} \leq \frac{m}{y}.$$

   Integrating from $y$ to $\hat{y}$ we get $$\hat{\theta} - \theta(y) \leq m \log(\hat{y}) - m \log(y).$$ Letting $y \to 0$ we get $$\hat{\theta} - \theta_\ast \leq - \infty,$$ which is an obvious contradiction. Therefore, it must be that $\theta_\ast = \pi/2$.
\end{proof}

We now prove the classification theorem for conformal solitons of the ICSF in $\h^2$.

\begin{theorem}
\label{t67}
   Let $\gamma$  be a conformal soliton of the ICSF solving \eqref{system-conformal} with initial conditions $(x(0),y(0),\theta(0))=(0,y_0,\pi/2)$. Then  $\gamma$ is a concave graph on the $x$-axis  over some bounded interval $(x_m, x_M)$, where the limits of the tangents at the endpoints are vertical.
    \end{theorem}

\begin{proof}
 From Proposition \ref{prg2}, we know that $\gamma$  is a graph on the $x$-axis. Let $y=y(x)$.   Due to Lemma \ref{lem-conformal-angles}, the angles between the tangent of $\gamma$ and the $x$-axis approach $\pm \pi/2$ at each of the ends. This implies that the domain of $y (x)$ is bounded, say,  of the form $(x_m, x_M)$ and that the tangents of $y(x)$ at the limit endpoints is $\infty$. To check the concavity of $y(x)$, the chain rule gives $y'(x)=\tan\theta$ and 
 $$y''(x)=\frac{\theta'}{\cos^3\theta}=-\frac{y+\cos^2\theta}{y\cos^4\theta}<0.$$
\end{proof}

%%%%%%
\subsection{Solutions of \eqref{system-conformal-reduced} with initial condition $\theta(0)> 0$.}

Consider  an orbit  $\alpha_\gamma = (y,\theta)$   of   \eqref{system-conformal-reduced} with the initial conditions $y(0) = y_0 > 0$ and $\theta(0) = \theta_0 > 0$. Then  $\theta'(0) = - \frac{1}{\cos\theta_0} - \frac{\cos\theta_0}{y_0} < 0$ and from Lemma \ref{le64}, $\theta(s)<0$. This implies that $y(s)$ is  increasing. Let $\lim_{s \to s_m} \alpha_\gamma(s) = (y_\ast, \theta_\ast)$. By   symmetry, we also have $\lim_{s \to s_M} \alpha_\gamma(s) = (y_\ast, - \theta_\ast)$.


We describe the conformal solitons in this case $\theta(0)$, proving that these conformal solitons fall in the family of Theorem \ref{t67}.  

\begin{theorem}
   Let $\gamma$  be a conformal soliton of the ICSF solving \eqref{system-conformal} with initial conditions $(x(0),y(0),\theta(0))=(0,y_0,\theta(0))$, with $\theta(0)>0$. Then  $\gamma$  can be reparametrized such that $\theta(0)=0$. In consequence, its geometric properties are described in Theorem \ref{t67}. 
 \end{theorem} 
\begin{proof}
It suffices to prove that there exists some $\overline{s} > 0$ such that $\theta(\overline{s}) = 0$.
    At $s = 0$ we have $\theta(0) = \theta_0 > 0$. Since $\lim_{s \to s_M} \theta(s) = - \theta_\ast < 0$, and using that $\theta$ is a decreasing function, there is   $\overline{s} > 0$ such that $\theta(\overline{s}) = 0$. Hence, $\alpha_\gamma$ can be considered as an orbit with initial condition $\theta(0) = 0$.  
\end{proof}

In  Fig. \ref{fig:conformal-solitons}, we show the plots of conformal solitons with varying initial conditions $\gamma(0)$.

\begin{figure}[h!t]
 \begin{center}
 \includegraphics[width=.4\textwidth]{conformal41.pdf}\qquad\qquad \includegraphics[width=.4\textwidth]{conformal42.pdf} 
\end{center}
\caption{Conformal solitons of the ICSF. In both cases, $y(0)=1$, with  $\theta(0)=0$ (left) and $\theta(0)=\pi/4$ (right).   }
\label{fig:conformal-solitons}
\end{figure}
 
%%%%%%%%%%%%%%%%%%%
 \section*{Ethics declarations}


The authors declare that they do not have any conflict of interest.

The authors have no relevant financial or non-financial interests to disclose.

No datasets were generated or analysed during the current study.



%%%%%%%%%%%%%%%%%%%%
 \section*{Acknowledgements}
Rafael L\'opez has been partially supported by MINECO/MICINN/FEDER grant no. PID2023-150727NB-I00, and by the ``Mar\'{\i}a de Maeztu'' Excellence Unit IMAG, reference CEX2020-001105- M, funded by MCINN/AEI/10.13039/ 501100011033/ CEX2020-001105-M.
%%%%%%%%%%%%%%%%%%%%%%%%
%%%%%%%%

\begin{thebibliography}{99}
%\bibliographystyle{plain}

%\bibitem{abcl}B. Andrews, B. Chow, C. Guenther, M. Langford. Extrinsic geometric flows. American mathematical society (206), 20202.


%\bibitem{al} U. Abresch, J. Langer, The normalized curve shortening flow and homothetic solutions. J. Differential Geom. 23, (1986), 175--196.

\bibitem{bda} B. D. Allen, Non-compact solutions to inverse mean curvature flow in hyperbolic space. Ph.D. thesis, University of Tennessee, Knoxville, 2016.

\bibitem{bl1} A. Bueno, R. L\'opez, Horo-shrinkers in the hyperbolic space. Taiwanese J. Math. 29 (2025), 1037--1059.

\bibitem{bl2} A. Bueno, R. L\'opez, The class of grim reapers in $H^2\times R$. J. Math. Anal. Appl. 541 (2025), 128730.

\bibitem{cl} I. Castro, A. M. Lerma,  Lagrangian homothetic solitons for the inverse mean curvature flow. Results Math. 71 (2017), 3--4.

\bibitem{cd} B. Choi, P. Daskalopoulos, Evolution of non-compact hypersurfaces by inverse mean curvature. Duke Math. J. 170 (2021), 2755--2803.

\bibitem{ch} B. Choi, P.K. Hung, Inverse mean curvature flow with singularities. Int. Math. Res. Not. IMRN 10 (2023), 8683--8702. 

%\bibitem{cz} K-S. Chou, X-P. Zhu, The curve shortening problem, Chapman-Hall/CRC, Boca Raton, FL, 2001.


%\bibitem{dt2} F. N. Da Silva, K. Tenenblat, Soliton solutions to the curve shortening flow on the 2-dimensional hyperbolic space. Rev. Mat. Iberoam. 38 (2022), 1763--1782.

\bibitem{dhw} G. Drugan, L. Hojoo, G. Wheeler, Solitons for the inverse mean curvature flow. Pacific J. Math.  284 (2016), 309--326.

%\bibitem{dt} H. F. S. Dos Reis, K. Tenenblat, Soliton solutions to the curve shortening flow on the sphere. Proc. Amer. Math. Soc. 147 (2019), 4955--4967.




%\bibitem{eg} C. L. Epstein, M. Gage, The curve shortening flow. In: Chorin A.J., Majda A.J. (eds) Wave Motion: Theory, Modelling, and Computation. Math. Sci. Res. Inst. Publ.  vol 7. Springer, (1987).

\bibitem{ga} M. E. Gage, Curve shortening makes convex curves circular. Invent. Math.  76 (1984), 357--364.

\bibitem{ge} C. Gerhardt, Flow of nonconvex hypersurfaces into spheres. J.   Differential Geom. 32 (1990), 299--314.

\bibitem{ge2} C. Gerhardt, Curvature problems. Ser. Geom. Topol. vol. 39. International Press, Boston. 2006.

\bibitem{ge3} C. Gerhardt, Inverse curvature flows in hyperbolic space. J. Differential Geom. 89 (2011), 487--527.



%\bibitem{gr} M. A. Grayson, Shortening embedded curves. Ann. of Math. (2), 129 (1989), 71--111.
 
%\bibitem{ha1} H. P. Halldorsson, Self-similar solutions to the curve shortening flow. Trans. Amer. Math. Soc. 364, (2012), 5285--5309.

%\bibitem{ha2} H. P. Halldorsson, Self-similar solutions to the mean curvature flow in the Minkowski plane $R^1_1$. J. Reine Angew. Math. 704 (2015), 209--243.
\bibitem{hi} G. Huisken, T. Ilmanen, The inverse mean curvature flow and the Riemannian Penrose inequality. J.    Differential Geom. 59 (2001), 353--437.

\bibitem{hw} P.-K. Hung, M.-T. Wang, Inverse mean curvature flows in the hyperbolic 3-space revisited. Calc. Var. Partial Differential Equations 54 (2015),   119--126.

\bibitem{kp} D. Kim, J. Pyo, Remarks on solitons for inverse mean curvature flow. Math. Nach. 293 (2020), 2363--2369.

\bibitem{lw} H. Li, Y.  Wei, On inverse mean curvature flow in Schwarzschild space and Kottler space. Calc. Var. Partial Differential Equations 56 (2017), no. 3, Paper No. 62.




 %\bibitem{lrm} R. F. de Lima, A. K. Ramos, J. P. dos Santos,  Solitons to mean curvature flow in the hyperbolic $3$-space. arXiv preprint arXiv:2307.14136, 2023.
 
 \bibitem{lu} S. Lu, Inverse curvature flow in anti-de Sitter-Schwarzschild manifold. Comm. Anal. Geom. 27 (2019), 465--489. 


 \bibitem{ma} L. Mari, J. R. Oliveira, A. Savas-Halilaj, R. Sodr\'e de Sena, Conformal solitons for the mean curvature flow in hyperbolic space. Ann. Glob. Anal. Geom. 65 (2024), 19. 
 
 \bibitem{ns} G. S. Neto, V. Silva,   Classification of ruled surfaces as homothetic self-similar solutions of the inverse mean curvature flow in the Lorentz-Minkowski 3-space. Bull. Braz. Math. Soc. New Ser. 54, Art. 48 (2023).
 
 \bibitem{pi1} G. Pipoli, Inverse mean curvature flow in quaternionic hyperbolic space. Atti Accad. Naz. Lincei Rend. Lincei Mat. Appl. 29 (2018), 153--171.

\bibitem{pi2} G. Pipoli,  Inverse mean curvature flow in complex hyperbolic space. Ann. Scient. \`{E}c. Norm. Sup. 52 (2019),  1107--1135.

\bibitem{sc} J. Scheuer, The inverse mean curvature flow in warped cylinders of non-positive radial curvature. Adv. Math.   306 (2017), 1130--1163.




%\bibitem{wx} E. Woolgar, R. Xie, Self-similar curve shortening flow in hyperbolic 2-space. Proc. Amer. Math. Soc. 150 (2022), 1301--1319.
 
 \bibitem{zh} H. Zhou,   Inverse mean curvature flows in warped product manifolds. J. Geom. Anal. 28 (2018), 1749--1772.



\end{thebibliography}
\
\end{document}

  




